\documentclass[11pt]{article}
\usepackage{fontspec}
\setmainfont{Times New Roman}
\setsansfont{Arial}
\setmonofont{Consolas}
\usepackage{amsmath,amssymb,mathtools}
\usepackage{geometry}
\usepackage{microtype}
\geometry{a4paper,margin=2.28cm}
\setlength{\parindent}{1.15em}
\setlength{\parskip}{0pt}
\makeatletter
\renewcommand\footnotesize{\@setfontsize\footnotesize{7.4}{8.4}\selectfont}
\makeatother
\providecommand{\mo}{\mathfrak{o}}
\providecommand{\ma}{\mathfrak{a}}
\providecommand{\ml}{\mathfrak{l}}
\providecommand{\mc}{\mathfrak{c}}
\providecommand{\mD}{\mathfrak{D}}
\emergencystretch=140em
\allowdisplaybreaks
\sloppy
\hbadness=10000
\vbadness=10000
\hfuzz=2pt
\vfuzz=10pt
\raggedbottom
\begin{document}
% Унит: Папер34 Сецтион20 соурце-фиделиты ремедиатион в001. Дирецт Интерславиц Латин ауториты транслатион фром the цуратед оригинал Герман сцан витнесс, нот фром the цомпрессед лоцал ТеКс.
\section*{34. Хиперкомплексне величины и теорија представјень}
\emph{Продолженье и завршенье: поглавје III, §§15--19, и поглавје IV, §§20--26.}

\section*{Поглавје IV. Представјеньа груп и хиперкомплексных системов}

\subsection*{§ 20. Уврстјенье хиперкомплексных системов в теорију}

Разматрајемо хиперкомплексне системы $\mo$ взходно к комутативному телу $P$. Понеже по договору за идеалы в $\mo$ в рачун приходет само $P$-модулы, за леве и праве идеалы важе условје максималности и условје минималности (ср. §8). Зато важи цела теорија колц, изложена в поглавју II.

Под представјеньами хиперкомплексного система $\mo$ в дальньем всегда разумејемо представјеньа в телу $P$, и то така представјеньа, при кторых из $a\sim A$ следује
\[
        a\rho\sim A\rho \qquad(\rho\in P),
\]
то јест модулна хомоморфност взходно к $P$. За модулы представјеньа то значи
\[
        (a\rho)m=a(m\rho)
\]
(ср. конец §15). Зато за модулы представјеньа важе теоремы §19; из ньих следује, же все иредуцибилне представјеньа посредујут се простыми левыми идеалами колца
\[
        \mo_0=\mo/\mc,
\]
и же зато имајемо толико нееквивалентных иредуцибилных представјень, колико двустранно простых суммандов $\ma_\nu$ појавјаје се в разложеньу
\[
        \mo_0=\ma_1+\cdots+\ma_s.
\]

Да бы се стварно определило представјенье од $\mo$, задано таким левым идеалом $\ml_\nu$, најпрво се преходи од елементов $\mo$ к соответным класам остатков в $\mo_0$. Понеже множенье с елементом из $P$ даваје изоморфизм за все идеалы в $\mo_0$, тело $P$ јест хомоморфно подтелу $e^{(\nu)}P$ тела автоморфизмов $K_\nu$ всакого простого левого идеала, и понеже јединичны елемент јест јединичны оператор, то јест чак изоморфизм. Представјенье в $P$, посредовано простым левым идеалом $\ml_\nu$, можно добити тако: најпрво изучити представјенье, посредовано чрез $\ml_\nu$ в том подтелу $e^{(\nu)}P$, а потом чрез изоморфизм $e^{(\nu)}P\simeq P$ прејдти к $P$. Тело автоморфизмов $K_\nu$ имаје конечны ранг взходно к $P$; зато иде о представјенье в подтелу конечного степена тела автоморфизмов од $\ml_\nu$.

Зато, да бы се добило искано представјенье, најпрво представимо всако $c$ из $\mo_0$ чрез јего двустранне компоненты:
\[
        c=c_1+\cdots+c_s.
\]
Тогда треба најдти само представјајучу матрицу од $c_\nu$, понеже друге $c_i$ анулирајут идеал $\ml_\nu$ и зато представјајут се нульевоју матрицоју. По §18 она се находи тако: представити $c_\nu$ чрез матричне јединице $\alpha_{ik}^{(\nu)}$ из $\ma_\nu$:
\[
        c_\nu=\sum c_{ik}^{(\nu)}\alpha_{ik}^{(\nu)}
\]
и в матрици $(\alpha_{ik}^{(\nu)})$ заменити всакы елемент $\alpha_{ik}^{(\nu)}$ матрицоју $A_{ik}^{(\nu)}$, ктора јему одговарја в представјеньу тела $K_\nu$ над $P$, посредованом самым $K_\nu$.

Приметка. Тело $K_\nu$ јест конечного ранга взходно к $P$. Всакы елемент $x$ задовольаје једначеньу $f(x)=0$ с коефициентами из $e^{(\nu)}P$, понеже меджу степеньами $x$ сигурно сушчествује линеарна зависност. Ако особливо $P$ јест алгебраично затворјено, то једначенье разпадаје се на линеарне факторы; понеже $K_\nu$ јест тело, $x$ уж јест нульево место једного линеарного фактора и зато принадлежи к $e^{(\nu)}P$. То јест $K_\nu=e^{(\nu)}P$. В том случају саме матрице $(\alpha_{ik}^{(\nu)})$ творет иредуцибилно представјенье.

Из тој приметкы заједно с концем §19 следује "теорема Бурнсида":

Нехај $P$ буде алгебраично затворјено и комутативно повезано с колцом $\mo$. Тогда всако иредуцибилно представјенье $n$-того степена над $P$ садржи точно $n^2$ линеарно независне матрице.\footnote{Зато $\mo$ не предполага се како хиперкомплексны систем и може напр. быти групово колцо, составјено из всих линеарных комбинациј по конечно многих елементов бесконечној групы.}

Из тој формы одмах следује обычна: "Систем матриц $n$-того степена над $P$, кторы заједно с всакоју двоју матриц садржи и јих продукт, јест иредуцибилны тогда и само тогда, когда не можно подати нетривиалну хомогенну линеарну релацију
\[
        \sum \alpha_{ik}x_{ik}=0
\]
с коефициентами из $P$, ктора јест задовольена елементами $x_{ik}$ всакој матрице система." Наиме, из всакого такого система матриц, додавши все линеарне комбинације, можно извести колцо и $P$-модул и сматрјати то колцо како јего властно представјенье.

Доказ. По §19 абсолутна област множительев $\mD$ јест изоморфна двустранно простому и пополно редуцибилному колцу $\mo_0$ с јединичным елементом; тако колцо јест полно матрично колцо (напр. колцо всих матриц $m$-того степена) взходно к телу автоморфизмов јего левых идеалов, кторо заради алгебраичного затворјеньа совпадаје с $P$. Всако иредуцибилно представјенье од $\mD$, зато и дано представјенье, поражджаје се простым левым идеалом. Левы идеал имаје ранг $m$, представјенье имаје степень $n$, зато $m=n$; тако в $\mo_0$, и зато такоже в $\mD$, сушчествује точно $n^2$ линеарно независне елементы.

Такоже лахко следује "обобчена теорема Бурнсида":\footnote{G. Frobenius und I. Schur, \emph{Über die Äquivalenz der Gruppen linearer Substitutionen}, Sitzungsber. Berlin 1906.}

Нехај под тыми самыми предположеньами $\mD$ буде пополно редуцибилно представјенье од $\mo$, кторо разпадаје се на саме нееквивалентне составне чести степеньев
\[
        n_1,n_2,\ldots,n_s;
\]
тогда $\mD$ имаје ранг
\[
        n_1^2+\cdots+n_s^2
\]
взходно к $P$.

Доказ. Нехај знову $\mo$ буде абсолутна област множительев, ктора необходно јест хиперкомплексны систем. Разне нееквивалентне представјеньа поражджајут се левыми идеалами $\ml_1,\ldots,\ml_s$, кторе принадлежет к разным двустранно простым компонентам $\ma_i$ од $\mo/\mc$, где $\mc$ јест радикал. Зато дано представјенье поражджаје се чрез $\ml_1+\cdots+\ml_s$ и имаје $\ma_1+\cdots+\ma_s$ како абсолутну област множительев. Из того, с помочу тых самых ранговых разсмотрень како выше, следује тврдженье.

Нехај ныне $\mo$ буде хиперкомплексны систем без радикала взходно к либовольному телу. По §13 $\mo$ јест пополно редуцибилно и имаје јединичны елемент. Из теорем §17 и §18 ныне следује: всако представјенье од $\mo$ јест пополно редуцибилно.

Обратно важи: всако пополно редуцибилно представјенье колца $\mo$, комутативно повезаного с $P$, имаје како абсолутну област множительев хиперкомплексны систем без радикала.

Доказ. Абсолутна област множительев $\mD$ јест изоморфна колцу и $P$-модулу из матриц над $P$, зато јест хиперкомплексны систем. Елементы радикала $\mc$ по §19 представјајут се нульеју в всаком иредуцибилном представјеньу, зато и в целом представјеньу; зато $\mc=0$.

(Ако в представјеньу наступајут само еквивалентне составне чести, кторе все сут посредоване једным левым идеалом $\ml$, тогда абсолутна област множительев ставаје двустранно проста.)

Како регуларно представјенье хиперкомплексного система $\mo$ означаје се представјенье, посредовано самым јединичным идеалом $\mo$. (При груповом колцу спецјално се выбира база из груповых елементов $u_1,\ldots,u_h$; ср. §6.) Понеже в $\mo$ како композицијне факторы сут присутне все леве идеалы од $\mo/\mc$, все иредуцибилне представјеньа уж находет се в регуларном представјеньу како диагоналне матрице $R_{ii}$, ако представити, же регуларно представјенье по формули (2) §16 чрез композицијны ред трансформује се к прикладној бази.

Представјенье јест знано, когда сут знане представјајуче матрице базовых елементов $u_1,\ldots,u_h$. Ако иде посебно о групово колцо, где $u_1,\ldots,u_h$ сут елементы групы, тогда всако хомоморфно матрично представјенье групы поражджаје такоже представјенье групового колца; зато проблем представјеньа групы јест специалны случај проблема представјеньа хиперкомплексных системов.

\end{document}

